\chapter{Wie wirkt Magie?}
\label{chap:magic}

\section{Was ist Magie?}

Magie gehört zu den außergewöhnlichsten Fähigkeiten, die ein Held in Heldenpfade erlernen kann. Magie kann Verletzungen heilen, Feuer beschwören, Verbündete stärken oder Gegner schwächen.

Nicht jeder Held kann zaubern. Ein Held benötigt das Talent \emph{Magiebegabt}, einen erlernten und verfügbaren Zauber sowie ausreichend Mana, um dessen Kosten zu bezahlen.

\section{Mana}

Jeder Zauber kostet Mana. Beim ersten Auftreten wird das Ausgeben von Mana als \textbf{Mana bezahlen (\manaLoss)} und das Erhalten von Mana als \textbf{Mana erhalten (\manaGain)} bezeichnet.

Jeder Held besitzt einen Manavorrat und aktuelles Mana. Der Manavorrat ist das Maximum, das aktuelle Mana steht für die momentan verfügbare Menge.

Zu Beginn jedes Abenteuers besitzt ein Held seinen vollständigen Manavorrat. Reicht das aktuelle Mana nicht aus, kann ein Zauber nicht gewirkt werden. Das aktuelle Mana kann niemals unter 0 oder über den Manavorrat steigen.

Verbrauchtes Mana wird während einer Rast mit \manaGain{} regeneriert.

\section{Magiebegabt}

Um Zauber wirken zu können, muss ein Held das Talent \emph{Magiebegabt} besitzen. Das Talent erhöht den Manavorrat des Helden um 6.

\section{Zauber und ihre Effekte}

Jeder Zauber wird durch eine Zauberkarte beschrieben. Die Karte nennt den Namen, die Manakosten, die Zauberzeit und den Effekt des Zaubers. Sie kann außerdem Ziele, eine Wirkungsdauer, Schlüsselwörter oder weitere Voraussetzungen angeben.

Ein \textbf{Zauber} ist die gesamte Karte. Um ihn zu wirken, bezahlt der Held seine Manakosten und führt den beschriebenen \textbf{Effekt} aus. Ein Effekt kann eine \textbf{Angriffsoption} sein. In diesem Fall bestimmt die Angriffsoption wie gewohnt die Angriffsfertigkeit, wie erzielte \star{} ausgegeben werden können, und mögliche zusätzliche Folgen. Jede aufgeführte Zeile darf pro Angriff höchstens einmal genutzt werden, sofern sie nicht ausdrücklich als mehrfach nutzbar bezeichnet ist; verbleibende \star{} können anschließend erhaltene Schadenswürfel verbessern.

Zauberkarten werden durch das Wirken eines Zaubers nicht automatisch erschöpft. Nur wenn eine Zauberkarte ausdrücklich das Symbol \textbf{Erschöpfen (\exhaust)} zeigt oder der Regeltext dies verlangt, wird genau diese Karte umgedreht. Danach gelten ausschließlich die Angaben und Effekte der nun sichtbaren Kartenseite.

Ein Held kann einen erlernten Zauber wiederholt wirken, solange genügend Mana vorhanden ist, die aktuell sichtbare Kartenseite das Wirken erlaubt und alle Voraussetzungen erfüllt sind.

\section{Einen Zauber wirken}

Jeder Zauber nennt eine Zauberzeit. Im Kampf gibt sie an, welche Aktion zum Wirken benötigt wird. Ein besonders schneller Zauber kann beispielsweise als \freeAction{} gewirkt werden, während ein aufwendigerer Zauber eine \standardAction{} benötigt.

\begin{kurzregel}[title=Einen Zauber wirken]
\begin{enumerate}
  \item Zauber wählen.
  \item Angegebene Manakosten mit \manaLoss{} bezahlen.
  \item Gültige Ziele bestimmen.
  \item Effekt vollständig ausführen.
  \item Falls angegeben, die Zauberkarte mit \exhaust{} erschöpfen.
\end{enumerate}
\end{kurzregel}

Kann ein Schritt nicht vollständig ausgeführt werden, kann der Zauber nicht gewirkt werden. Insbesondere müssen ausreichend Mana und mindestens ein gültiges Ziel vorhanden sein.

\begin{beispiel}[title=Beispiel: Einen unterstützenden Zauber wirken]
Lina möchte mit dem Zauber \emph{Heilender Funke} einen benachbarten Verbündeten heilen. Der Zauber kostet 1\manaLoss{} und benötigt eine \moveAction. Lina wählt den Zauber, bezahlt 1 Mana, bestimmt den benachbarten Verbündeten als gültiges Ziel und führt anschließend den Effekt aus. Der Verbündete erhält 4\lifeGain{} zurück.
\end{beispiel}

\section{Angriffszauber}

Enthält der Effekt eines Zaubers eine Angriffsoption, gelten für diese vollständig die Kampfregeln aus \chapterref{chap:combat}. Die Angriffsoption wird erst genutzt, nachdem die Manakosten des Zaubers bezahlt wurden.

\begin{kurzregel}[title=Angriffszauber]
Manakosten bezahlen $\rightarrow$ Angriffsoption nach den normalen Kampfregeln ausführen.
\end{kurzregel}

Bestimmt ein Angriffszauber nach der Angriffsprobe weitere Ziele und überträgt alle gekauften Schadenswürfel und Effekte auf diese, werden die Ziele vor dem Schadenswurf gewählt. Der gemeinsame Schadenspool wird nur einmal gewürfelt und das Ergebnis gegen jedes Ziel angewendet. Schilde und zielabhängige Boni werden für jedes Ziel einzeln berücksichtigt.

\begin{beispiel}[title=Beispiel: Eisblitz]
Lina wirkt \emph{Eisblitz}. Der Zauber kostet 1\manaLoss{} und benötigt eine \standardAction. Sein Effekt ist eine Angriffsoption mit der Angriffsfertigkeit \emph{Zaubern}. Bei dieser Angriffsoption kann 1\star{} für 1\redDmgDie{} oder 2\star{} für 2\redDmgDie{} ausgegeben werden. Für 3\star{} kann das Ziel stattdessen den Statuseffekt \emph{Festgefroren} erhalten.

Lina bezahlt 1 Mana, bestimmt einen gültigen Gegner als Ziel und legt anschließend die Angriffsprobe mit \emph{Zaubern} ab. Sie erzielt 2\star{} und gibt sie für 2\redDmgDie{} aus. Der weitere Ablauf folgt den normalen Regeln für Schaden, Rüstungsbruch, \armor{} und Lebenspunktverlust.
\end{beispiel}

\section{Zauberproben}

Nicht jeder Zauber verlangt eine Probe. Verlangt ein Effekt eine Probe, nennt die Karte Attribute, Fertigkeit und Schwierigkeit. Misslingt die Probe, tritt die Wirkung nicht ein; das ausgegebene Mana bleibt verbraucht.

Eine Angriffsprobe durch eine Angriffsoption ist ebenfalls eine Zauberprobe. Für sie gelten jedoch die Regeln für Angriffe aus \chapterref{chap:combat}.

\section{Ziele und Wirkungsdauer}

Die Zauberkarte bestimmt mögliche Ziele. Typische Zielangaben sind Selbst, Berührung, ein Ziel, mehrere Ziele oder ein Bereich. Ein Ziel muss grundsätzlich wahrnehmbar sein.

Die Reichweite einer Angriffsoption ergibt sich aus ihrer Angriffsfertigkeit, sofern der Zauber nichts anderes bestimmt.

Typische Wirkungsdauern sind sofort, bis zum Ende des Zuges, bis zum Ende der Runde oder eine bestimmte Dauer. Eine Kampfrunde endet, nachdem alle Gegner vollständig aktiviert wurden.

Endet ein Zauber, enden auch alle durch ihn verursachten fortdauernden Boni, Mali und Zustände, sofern die Karte nichts anderes bestimmt.

\section{Ausführliches Beispiel}

Lina möchte im Kampf zwei entfernte Gegner mit \emph{Eisblitz} angreifen.

\begin{enumerate}
  \item \textbf{Zauber wählen:} Lina wählt den verfügbaren Zauber \emph{Eisblitz}. Das Wirken benötigt eine \standardAction.
  \item \textbf{Mana bezahlen:} Sie bezahlt die Grundkosten von 1\manaLoss.
  \item \textbf{Ziel bestimmen:} Da die Angriffsoption \emph{Zaubern} verwendet, darf sie einen benachbarten oder entfernten Gegner als erstes Ziel wählen.
  \item \textbf{Angriffsprobe:} Lina legt eine Angriffsprobe mit \emph{Zaubern} ab und erzielt 3\star.
  \item \textbf{Sterne ausgeben:} Lina gibt 1\star{} für 1\redDmgDie{} und 2\star{} für 2 weitere \redDmgDie{} aus. Ihr Schadenspool besteht aus 3\redDmgDie{}.
  \item \textbf{Weiteres Ziel wählen:} Lina bezahlt zusätzlich 1\manaLoss{} und wählt einen zweiten entfernten Gegner in Reichweite. Beide Ziele sind von allen gekauften Schadenswürfeln betroffen.
  \item \textbf{Angriff abschließen:} Lina würfelt den gemeinsamen Schadenspool einmal. Das Ergebnis wird gegen beide Gegner angewendet; \armor{} und mögliche weitere Effekte werden für jedes Ziel einzeln berücksichtigt.
\end{enumerate}

\section{Weiterführende Informationen}

Die vollständige Zauberliste befindet sich in \appendixref{app:spells}.
